% Pursuing Stacks (cotes 134-2 to 134-8): archive cote 134-2, site page 21 (PDF page 22).
% letter to Quillen. Draft compared: ps1.tex:306-389.
% AI-transcribed by Claude Opus 5.5 (claude-opus-5-5, adaptive thinking, effort high), Message Batches API, batch ps1 (msgbatch_01Rv5hCaNTtUt4JHx1yqoHPL), 2026-11-12; not yet checked by eye.
\scanpage{134-2/21}
\mspage{\hand{4}}

yesterday.

What may look less conviencing though at first sight, is my choice
to view \inserted{as non primitive} even the "degeneration maps" $k^{i}_{1} : F_{i} \to F_{i+1}$ , associating to
every i-object the \hand{(}i+1\unsure{\hand{)}}-object acting as an identity on the former, In
all cases I have met so far, these maps are either given beforehand with
the structure (of a 1-category or 2-category\hand{,} say), or they can be uniquely
deduced from the axioms. In the present set-up however, they seem to me
to appear more naturally as ``primary'' (not as primitive) structure, much
in the same way as the $\mathrm{inv}_{i}$ . Different choices for associating an $\infty$ -
groupoid to a topological space, while yielding the same base-sets $F_{i}$ ,
will however (according to this point of view) give rise to different
maps $k^{i}_{1}$ . The main motivation for this point of view comes from the fact
that the mechanism for a uniform construction of the chain of ever higher
order structures makes a basic use of the source and target maps only and
of their ``transposés'' (see below), and (it seems to me) not at all of the
degeneration maps, which in this respect \inserted{are} rather \deleted{are} confusing \deleted{\unsure{x}} \inserted{the real} picture \unsure{\over{.}{,}}
\inserted{if viewed as "primitive".}
\margin{(}The degeneration maps rather appear as typical cases of primary structure,
probably of \deleted{\unsure{a}}special significance in the practical handling of $\infty$ -grou-%
poids, but\hand{/}not at all in the conceptual machinery leading up to the cons\xout{\illegible{1}}
struction of \unsure{the|thu} structure species \deleted{\illegible{1}}\hand{``}$\infty$ -groupoids\hand{''}.

\margin{$\overset{i}{\underset{1}{*}}$}
Much in the same way, the composition operations $\overset{i}{\underset{1}{*}}$ are viewed as
primary, notas primitive or skeletal structure. Their description for
the fundamental $\infty$-groupoid of a space - for instance the description of
composition of paths - depends on arbitray choices, such as the choice of
an isomorphism (say) between $(I,1)\amalg_{e}(I,0)$ and I, where I is the unit
interval, much in the same way as the notion of an inverse of a path
depends on the choice of an isomorphism of I with itself, exchanging the
\margin{$\overset{i}{\underset{2}{*}}$}
two end-points 0 and 1. The operations $\overset{i}{\underset{2}{*}}$ of Godement take sense only once
the composition operations $\underset{1}{*}$ are defined - they are ``seconndary structu-%
re'' , and successively the operations $\underset{3}{*}$ \ldots $\underset{i}{*}$ appear as ternary etc
structure. This is correctly suggested by the notations which I chose\xout{\illegible{1}}
yesterday, \unsure{\over{w}{W}}here however I hastily \deleted{\unsure{throwed}}\inserted{threw} all these operations into a same
pot baptized ``primary structure''\unsure{\hand{!}}

\over{4}{5}. It is about time though to come to a tentative precise definition of
description of the process of stepwise introduction of an increasing chain
of higher order structure. This will be done by intoducing a canonical
sequence of categories and functors
\[ C_{0} \to C_{1} \to C_{2} \to \ldots \to C_{n} \to C_{n+1} \to \ , \]
where $C_{n}$ denotes the \xout{xxxxxxxxxx} category harbouring the ``universal'' parti-%

\mspage{\hand{21}}
